\chapter {\nameTitleSpace LITERATURE REVIEW}
\chpterTitleSection
\section{Conceptual Framework}
\sectionTitleBody
\subsection{Concept of Plant Disease Detection}
\sectionTitleBody
Farmers face major challenges in dealing with crop disease,
	making classification and examination of illnesses key to optimizing
	food yield in agriculture \citep{chenIdentificationFruitTree2021a}. The study of detecting
	and identifying plant disease is important, especially as it can potentially
	monitor extensive crop fields and rapidly identify disease symptoms
	on plant leaves. Therefore, the desire for a quick, accurate and cost
	effective approach to predicting instances of crop disease is of utmost
	importance \citep{chenIdentificationFruitTree2021a}. Artificial Intelligence (AI) greatly
	contributes to the agricultural sector, improving a nation\textquoteright s
	gross domestic product (GDP). The application of edge intelligence
	to agriculture, especially through the use of deep learning models
	like the YOLOv3 neural network on embedded systems such as the NVIDIA
	Jetson TX2, presents a novel approach. Hyperspectral and multispectral
	knowledge is widely employed in different application areas in agriculture,
	including sustainable agriculture and these technologies provides necessary
	knowledge for farmers and agricultural management \citep{angBigDataMachine2021a}


\subsection{Traditional Methods of Disease Detection} \sectionTitleBody
Traditional plant disease detection is through manual observation
	by farmers and other agricultural extension workers. However, in majority
	of the cases, the diagnosis is based on guess on the disease which
	is performed visually by humans. Trained raters may be efficient in
	recognizing and quantifying diseases but they have associated limitations
	that may cause harm to crop yield. Challenges of this method are listed
	below: 
\begin{enumerate}
	\item Accuracy depends on the experience and accuracy of the observer.

	\item Diseases are diagnosed only when symptoms are visible, preventing
	early intervention. 

	\item It is difficult to differentiate between diseases
	with similar symptoms by mere visual observation. 
	
	\item  Manual inspection is susceptible to error. 
\end{enumerate}

\subsection{Overview of Artificial Intelligence and Machine Learning  in Agriculture}

\subsubsection{Deep Learning in Agriculture}
\sectionTitleBody
Disease prediction at an early stage is vital to preventing
	massive crop loss and ensuring greater crop yield. Advancement in
	technology in terms of speed and power, along with improved access
	to massive datasets, contribute to the enhanced robustness of disease
	detection systems \citep{wuDetectionEnumerationWheat2020a}. A hybrid integration of classifier
	methods showed a recognition rate of 91.11 \%, surpassing the performance
	of serial, parallel, and deep learning approaches \citep{elmassiCombinationMultipleClassifiers2021a}.
	In addition, the application of the CovNet algorithm for recognizing
	weeds in crops and the use of hybrid deep-learning models demonstrate
	potential results in terms of accuracy and parameter reduction \citep{zhouTomatoLeafDisease2021a}. Deep transfer learning, leveraging pre-trained datasets
	like Inception and ImageNet modules, proves effective in identifying
	plant diseases, with the aid of Support Vector Machine (SVM) and multi-layer
	perceptron contributing to higher accuracy. The deployment of Deep
	Convolutional Neural Networks for identifying corn dietary sickness
	and the use of hardware enhancements, such as Raspberry Pi3 with an
	Intel Movidius Neural Compute Stick, has lead to superior metric accuracy
	and performance \citep{sunNorthernMaizeLeaf2020a}.

As AI and DL technology advances, Computer Vision (CV) has
	made significant strides. CV-based approaches employ Principle Component
	Analysis (PCA) that has proven useful in diagnosing grape leaf diseases
	with quality research accuracy \citep{xieDeepLearningBasedRealTimeDetector2020a}. Validation accuracy
	leveraging VGGNet on real-world datasets reached 91.83\% \citep{xieDeepLearningBasedRealTimeDetector2020a}. Vision Transformer (ViT) has shown potential in developing
	models for detecting plant diseases. Researchers have examined ViT
	to identify plant disease on mobile devices and classify plant automatically
	in real-time. ViT based models out performed prior existing model
	in classifying benchmark datasets such as CIFAR-100, Oxford-IIIT pets,
	Oxford Flowers-102, and ImageNet.

\subsubsection{Image Based Techniques} \sectionTitleBody
Image - based disease detection has become powerful and reliable
	in modern precision agriculture. This method utilizes digital images
	of plant leaves and advanced computational algorithms to identify
	diseases based on features such as colour, textual and shape. With
	the integration of machine learning and deep learning approach, particularly
	the CNNs, image based systems can detect and classify disease with
	high precision. The existing system supports two popular machine learning
	models, Support Vector Machine (SVM) and Deep Learning. Based on
	the result suggestions were preferred for image processing based plant
	leaf disease detection within the constraints of architectural complexity,
	computational power, and training data. Stages are shown in figure~\ref{fig:image_processing_workflow} below:

\begin{figure}[!h]
	\centering\includegraphics[height=\linewidth]{image_processing_workflow.png}
	\caption{Image Based Techniques Steps \citep{bansalDiseaseDetectionApple2021} \label{fig:image_processing_workflow}}
\end{figure}


 %\subsubsection{Studies on Ugu Leaf Disease Detection}The survey of leaf spot disease of fluted pumpkin was carried out in 2019 during farming season. Leaf samples were collected from individual farmers at Ten Telfairia locations where Telfairia were cultivated namely: Apir, Federal Housing Estate, Kanshio, Ichuwa, Welfare quarters, New GRA, Yarkyo, and Lower Benue River Basin, along River Benue Bank and at NASME (Nigeria Army School of Military Engineering) barracks. Disease incidence was examined by expressing the numbers of infected Telfairia plants as a percentage of the total number of Telfairia plants.     (Ekhuemelo \& Otor, 2021).

\newpage
\subsubsection{Studies on Pumpkin Leaf Disease Detection} \sectionTitleBody
The survey of leaf spot disease of fluted pumpkin was carried
	out in 2019 during farming season. Leaf samples were collected from
	individual farmers at Ten Telfairia locations where Telfairia were
	cultivated namely: Apir, Federal Housing Estate, Kanshio, Ichuwa,
	Welfare quarters, New GRA, Yarkyo, and Lower Benue River Basin, along
	River Benue Bank and at NASME (Nigeria Army School of Military Engineering)
	barracks. Disease incidence was examined by expressing the numbers
	of infected Telfairia plants as a percentage of the total number of
	Telfairia plants. \citep{ekhuemeloASSESSMENTFUNGIASSOCIATED2021}.

\newpage
\section{Theoretical Framework} \nameTitleSpace
The foundation of this study is built on the theories below:
\subsection{Decision Theory in Artificial Intelligence}\sectionTitleBody
Decision theory is a foundational concept in Artificial Intelligence
	(AI), which enables machines to make analytical and informed decisions
	based on available data. It gathers principles from various fields
	such as mathematics, statistics, economic and psychology to model
	and improve decision making processes. Decision theory provides a
	platform to predict outcomes, evaluate choices and guide actions in
	uncertain environment. Applying this theory in the study will enable
	AI and ML model to imitate expert diagnosis in recognizing pattern
	and proffering treatment based on disease type detected.

\subsection{Pattern Recognition Theory} \sectionTitleBody
Pattern Recognition is the process of using machine learning
	algorithms to recognize patterns. It means sorting data into categories
	by analyzing the pattern present in the data. One major benefit of
	pattern recognition is that it can be used in many areas. In pattern recognition
	application, raw data is processed and converted to the form the machine
	can use. The two major method involved in pattern recognition is clustering
	and classification. Pattern recognition will allow for extraction
	of disease feature in pumpkin leaf images and classify diseases based
	on learned visual variances.

\section{Empirical Framework} \nameTitleSpace
\subsection{Related Work}
\sectionTitleBody
Plant disease detection through machine learning
	and deep learning models have been an interesting area of recent research. \cite{jiangRealTimeDetectionApple2019a} developed a deep CNN architecture for identifying
	apple leaf disease. The architecture integrates CNN with the inception
	module to recognize five types of apple leaf disease. \cite{bukumiraCarrotGradingSystem2022a} proposed a carrot grading system using computer vision
	techniques and a CNN graph. These identified frameworks showed the capability
	to classify carrot fruit accurately and grade them. In addition, other
	studies have suggested DL models for recognizing apple leaf diseases
	including Ful et al. (2022), \cite{bansalDiseaseDetectionApple2021} and \cite{khanDeepDiagnosisRealtime2022a} To enhance accuracy, the models employ different methods like
	image augmentation, attention mechanisms, and dilated convolution.
	Studies have examined hybrid deep-learning models in diverse fields,
	such as agriculture and medical sector, for disease diagnosis and
	classification \citep{zhangInstancebasedDeepTransfer2023a} and \citep{waleedNovelDeepLearning2022}. These
	investigations have proved the efficacy of combining deep learning
	architectures to accurately extract relevant features from the input
	data. \cite{hosnyMultiClassClassificationPlant2023a} designed a hybrid approach using lightweight
	deep model features and Local Binary Pattern features to detect and
	classify leaf diseae in grapes, apples and tomatoes. These techniques
	resulted in accuracy rate of 97.8\%, 97\% and 96\% for each types
	of plants respectively. \citep{ahilaMaizeLeafDisease2019} create CNN model
	by adjusting the LeNet model to classify corn leaf disease. A corm
	leaf image dataset from PlantVillage was utilized to train their method.
	The techniques resulted in an accuracy of 97.89\%. \citep{waheedOptimizedDenseConvolutional2020a}
	developed an optimized DenseNet structure to detect and classify disease
	which includes northern leaf blight, common rust, gray leaf spot,
	and healthy leaves leading to an accuracy of 98.06\%. Wu et al. (2019)
	proposed a CNN model to detect northern leaf blight on corn leaves
	using UAV images. They used the pre-trained ResNet architecture for
	feature extraction, employing its result as inputs to a basic linear
	classifier leading to an overall accuracy of 95.1.
\newpage
\section{Summary of Related Works for Plant Leaf Disease Detection and Classification}

\begin{center}
% \begin{longtable}[!h]
	% \caption{\label{tab:key_ai} Summary of Related Works for Plant Leaf Disease Detection and Classification}
	%{|p{5cm}|p{9cm}|}
	\begin{longtable}{|c|p{2cm}|p{2cm}|p{2cm}|p{2cm}|}
		\caption[Related Works Summary]{Summary of Related Works for Plant Leaf Disease Detection and Classification}
		\label{tab:key_ai} \\
		%%%%%%%%%%%%%%%%%%%%%%%%%%%%%%%%%%%%
	
		\endfirsthead
		
		\multicolumn{2}{c}%
		{{\bfseries \tablename\ \thetable{} -- continued from previous page}} \\
		\hline
		\bfseries References	& \bfseries Objective	& \bfseries Method & \bfseries Key Finding & \bfseries Limitations\\
		\hline
		\endhead
		%%%%%%%%%%%%%%%%%%%%%%%%%%%%%%%%%%%%
		\hline
		\bfseries References	& \bfseries Objective	& \bfseries Method & \bfseries Key Finding & \bfseries Limitations \\
		
		\hline
		\cite{hosnyMultiClassClassificationPlant2023a}	& Multi-class plant disease classification	& Hybrid of CNN and Local Binary Pattern	& Improved multi-class classification accuracy &	Complexity in feature integration may stop model interpretability and adaptability.\\
		
		\hline 
		\cite{bukumiraCarrotGradingSystem2022a} & Graph Convolutional Neural Network (GCNN) and computer vision on carrot grading & Cascaded graph CNN and Bayesian Optimisation & Automated grading of carrot & Limited carrot grading only. The GCNN may be computationally intensive for real-time grading\\

		\hline 
		Fu et al. (2022) & Lightweight model for apple leaf disease detection &	Lightweight CNN	& Reduced model complexity with increased accuracy &	Limited to apple leave, may lack reliability for other crops.\\
		
		\hline
		\cite{khanDeepDiagnosisRealtime2022a} &	Real-time apple detection system &	Lightweight-CNN and YOLOv4 &	Real-time performance &	Classification performance requires improvement.\\
		\hline
		\cite{bansalDiseaseDetectionApple2021}	& Apple leaf disease detection &	Ensemble of pre-trained deep CNN models.	& Accuracy and enhanced classification for apple disease. &	Multiple diseases may reduce accuracy.\\
		\hline
		\cite{waheedOptimizedDenseConvolutional2020a} &	Disease recognition and classification of corn leaf &	Optimized DenseNet Model&	Enhanced classification effectiveness and accuracy &	The dense architectural structure may increase memory consumption.\\	
			
		\hline
		\cite{ahilaMaizeLeafDisease2019}
		 &	Classification of maize leaf diseases &	Binary pattern modified LeNet &	Real-time performance with Lightweight model &	Limited to maize leaf disease classification.\\
		
		\hline
		\cite{jiangRealTimeDetectionApple2019a} & Real-time detection of disease on apple leaf & Enhanced CNN & Provide high quality performance solution with real time detection & Complexity may result to challenge in low-power devices.\\
		
		\hline
		Wu et al (2019) 	&Autonomous detection of plant disease&	CNN model on aerial imagery	& Robust in large-scale disease detection&	Limited to aerial imagery.\\
		\hline
		
	\end{longtable}
% \end{longtable}
\end{center}

\subsection{Limitations of Existing System}\sectionTitleBody
\begin{enumerate}
	\item Complexity may result to challenge in low powered devices.
	\item May lack reliability for other crops apart for the preset
	ones.
	\item Classification performance requires improvement.
	\item Complexity in feature integration may limit model easy of
	interpretation.
	\item The dense architectural structure may increase memory consumption.
\end{enumerate}

\subsection{Research Gaps} \sectionTitleBody
% From table~\ref{tab:key_ai}, \citep{hosnyMultiClassClassificationPlant2023a} built a multi-class plant disease classification system using hybrid CNN and Local Binary Patter, producing a very complex sytem which is difficult to adapt. \cite{bukumiraCarrotGradingSystem2022a} used a Graph Convolutional Neural Network (GCNN) and computer vision, and cascaded graph CNN and Bayesian Optimisation to build a carrot grading model, resulting in a computationally intensive system for real - time appication. \cite{khalidRealTimePlantHealth2023a} built a real-tie apple detection system using Lightweight - CNN and YOLOv4 which produce a system with real-time performance that needs performance improvement. \cite{bansalDiseaseDetectionApple2021} got an apple leaf disease detection with ensemble pre-trained deep CNN model for multiple diseases but with reduced accuracy, even for apple leaves only. 
%\cite{waheedOptimizedDenseConvolutional2020a} built a disease recognition and classification of corn leaf disease using optimised densenet model, the resultant model was was intense on memory, increasing the cost of affordability. While \cite{ahilapriyadharshiniMaizeLeafDisease2019a} came with a maize classification leaf disease using binary pattern modified LeNet, which is lightweight and real-time but limited to corn / maize. \cite{jiangRealTimeDetectionApple2019a} was real-time detection of disease on apple leaf with enhanced CNN providing high quality performance solution  with real time detection but with low power challenge for low power devices like mobile phones, tablets and IoT devices. With (Wu et al. 2019) was an Autonomous detection of plant disease built with CNN model on aerial imagery, robusst in large scale diseae detection but limited to aerial imagery, hence not easily adaptable to everyday farmers.\\

% From the existing systems above analysed, there is a gap of a real - time lightweight mobile model deployable on the low powered mobile systems like smart phones, tables, raspbery pi and raspbery pi pico. This is the gap that this work is going to fill. 

The related work shows gap in research investigating only one disease type on a particular leaf plant. The existing system can detect only one type of disease on pumpkin leaf. And this project is to fill this gap by covering four different diseases of the pumpkin leaf.






